% Generated draft: requires mathematical and source review.
\documentclass[11pt]{article}
\usepackage[margin=25mm]{geometry}
\usepackage{fontspec}
\setmainfont{Latin Modern Roman}
\usepackage{amsmath,amssymb,mathtools,unicode-math}
\setmathfont{Latin Modern Math}
\usepackage{enumitem,needspace,xurl,hyperref}
\setlength{\parindent}{0pt}
\setlength{\parskip}{5pt}
\setlength{\emergencystretch}{3em}
\newcommand{\N}{\mathbb{N}}
\newcommand{\Z}{\mathbb{Z}}
\newcommand{\Q}{\mathbb{Q}}
\newcommand{\R}{\mathbb{R}}
\newcommand{\C}{\mathbb{C}}
\newcommand{\F}{\mathbb{F}}
\newcommand{\A}{\mathbb{A}}
\newcommand{\PP}{\mathbb{P}}
\newcommand{\E}{\mathbb{E}}
\newcommand{\eps}{\varepsilon}
\newcommand{\dd}{\,\mathrm{d}}
\begin{document}
\section*{UnsolvedMath ID: 30000384}
\textbf{Source problem code:} catalogue entry
\subsection{Definitions and mathematical statement}
For a surface \(S\) in the intended class of geometric polyhedral surfaces, write \(f_0(S)=n\) for its number of vertices and \(f_1(S)\) for its number of edges. The source states a quadratic upper bound \(f_1(S)=O(n^2)\) and asks whether that order is asymptotically tight.

A precise lower-bound interpretation is the existence of a constant \(c>0\) and a sequence of admissible surfaces \(S_j\) with
\[
 n_j=f_0(S_j)\longrightarrow\infty,\qquad
 f_1(S_j)\geq c n_j^2\quad\text{for every }j.
\]
The constant must be independent of \(j\), and the surfaces must satisfy the same geometric realization requirements as those in the stated upper bound. A stronger all-large-\(n\) formulation would require such a surface for every sufficiently large exact vertex count, not just an unbounded sequence. The source does not distinguish these versions of tightness.

In a common convention, a geometric polyhedral surface is a two-dimensional manifold realized by straight edges and planar polygonal faces meeting along common faces, often embedded without self-intersection in \(\R^3\). The supplied record does not specify the ambient dimension, embeddedness, closedness, connectedness, orientability, allowance of boundary, or whether faces must be triangles. That conventional setting is explanatory only, not a silently imposed class. Abstract combinatorial surfaces and immersed or intersecting realizations cannot be substituted for geometric realizations without confirmation.

In the usual embedded-polyhedral-complex convention the edge graph is simple, giving \(f_1\leq\binom n2\), which explains a quadratic combinatorial ceiling without supplying geometric surfaces that approach it. Restricting to convex surfaces or fixed topology would change the problem. For example, for closed orientable triangulated surfaces of fixed genus \(g\), Euler's relation gives \(f_1=3n-6+6g\); quadratic edge growth in that setting would require the topology to vary with \(n\). No fixed-genus or convexity restriction is added to the source's question.

The literature triage checked 2026-08-21 reports that tightness of the geometric quadratic bound is not known and that published constructions remain far below the combinatorial extremal order. It gives no quantitative lower-bound rate. This assessment and the verification field's source-checking report remain dated metadata rather than independently verified current claims.
\subsection{Short English statement}
Can admissible geometric polyhedral surfaces have at least a fixed positive multiple of the square of their vertex count in edges, for unbounded vertex counts? The geometric class and meaning of asymptotic tightness need confirmation. A combinatorial quadratic ceiling alone does not construct the required geometric surfaces.
\subsection{Sources}
\begin{itemize}
\item Local UnsolvedMath registry, permanent ID 30000384, source code catalogue entry, ``Maximum Edge Counts of Geometric Polyhedral Surfaces,'' including original and cleaned asymptotic question and verification metadata.
\item \emph{Combinatorics (2006)}, Oberwolfach Reports, pp.\ 5--72, DOI 10.4171/owr/2006/01: \url{https://doi.org/10.4171/owr/2006/01}.
\item Supplied reference ``polyhedral surfaces of high genus'': \url{https://arxiv.org/abs/math/0412093}.
\item Literature triage checked 2026-08-21; local record updated 2026-09-16. No external source has been visited or independently checked here.
\end{itemize}
\end{document}
